% Constitutional note:
% Chapter 34 (concordance.tex) is the normative status ledger.
% If a local statement disagrees with Chapter 34's status or scope,
% the local statement is stale and must be rewritten or ignored.
\chapter{Introduction}

\label{rem:volume-one-governing-question}

\section{The logarithmic 1-form}
\label{sec:the-seed}
\index{propagator|textbf}

\index{Koszul duality!classical}
Consider two distinct points $z_1, z_2$ on a smooth algebraic
curve~$X$, and the logarithmic 1-form
\begin{equation}\label{eq:the-seed}
\eta_{12} \;=\; d\log(z_1 - z_2)
\;=\; \frac{dz_1 - dz_2}{z_1 - z_2}
\end{equation}
on the configuration space $C_2(X) = \{(z_1,z_2) \in X^2 : z_1 \neq z_2\}$.
It has a simple pole along the diagonal $z_1 = z_2$ with
$\operatorname{Res}_{z_1 = z_2} \eta_{12} = 1$, and it is
conformally invariant.

For three distinct points $z_1, z_2, z_3 \in X$, the three such
forms satisfy the Arnold relation
\begin{equation}\label{eq:arnold-seed}
\eta_{12} \wedge \eta_{23}
+ \eta_{23} \wedge \eta_{31}
+ \eta_{31} \wedge \eta_{12} = 0,
\end{equation}
which follows from the partial fraction identity: the 2-form
$\eta_{12} \wedge \eta_{23}$ has poles along $D_{12}$ and $D_{23}$,
and the identity~\eqref{eq:arnold-seed} is the assertion that
the signed sum of residues along the three pairwise collision
divisors of $\overline{C}_3(X)$ vanishes.

Now let $\cA$ be a chiral algebra on~$X$.  Place $n$ of its elements
$a_1, \ldots, a_n$ at positions $z_1, \ldots, z_n$ and form
\begin{equation}\label{eq:bar-element-intro}
a_1(z_1) \otimes \cdots \otimes a_n(z_n)
\otimes \bigwedge_{(i,j) \in E(T)} \eta_{ij}
\;\in\;
\cA^{\boxtimes n} \otimes
\Omega^{n-1}_{\log}(\overline{C}_n(X)),
\end{equation}
summed over edges of a planar tree~$T$.  Define a differential $d$
by extracting residues along collision divisors:
$d(a_1 \otimes \cdots \otimes a_n) = \sum_{i < j}
\operatorname{Res}_{D_{ij}}[\eta_{ij} \cdot
\text{OPE}(a_i, a_j)] \otimes (\text{remaining factors})$.
The Arnold relation~\eqref{eq:arnold-seed} forces $d^2 = 0$:
\[
d^2(\cdots)
= \sum_{i < j < k}
\operatorname{Res}_{D_{ij}} \operatorname{Res}_{D_{jk}}
[\eta_{ij} \wedge \eta_{jk} \cdot (\cdots)]
+ \text{cyclic permutations}
= 0,
\]
because~\eqref{eq:arnold-seed} makes the triple residue vanish.
The resulting cochain complex is the \emph{bar complex}
$\barB_X(\cA)$.

\begin{remark}[From the bar complex to its dual]
The cohomology $H^*(\barB_X(\cA))$ is the Koszul dual
$\cA^! = T(V^*)/R^\perp$.  Verdier duality on
$\operatorname{Ran}(X)$ intertwines the bar of $\cA$ with the bar of
$\cA^!$ (Theorem~\ref{thm:bar-cobar-isomorphism-main}):
\begin{equation}\label{eq:NAP-intro}
\int_X \cA \;\simeq\;
\left(\int_{-X} \cA^!\right)^{\!\vee}.
\end{equation}
The logarithmic form~$\eta_{ij}$ is the kernel of a non-abelian integral
transform; Verdier duality is its inversion formula; the cobar
construction $\Omega(\barB(\cA))$ reconstructs $\cA$ from the
transformed data.
\end{remark}

\begin{remark}[The Fourier perspective]\label{rem:intro-fourier-perspective}
The bar construction is a one-dimensional non-abelian Fourier transform
whose kernel is the logarithmic propagator~$\eta_{ij}$.  Its four
structural properties---intertwining, inversion, Plancherel,
characteristic function---are Theorems~A--D.  The explicit computations
of Chapter~\ref{ch:fourier-seed} exhibit this identification for the
Heisenberg and Kac--Moody families before the general theory.
\end{remark}

\begin{remark}[The nilpotence-periodicity correspondence]
\label{rem:nilpotence-periodicity-intro}
\index{nilpotence-periodicity correspondence}
The logarithmic form~$\eta_{12} = d\log(z_1 - z_2)$ is not an
incidental choice of kernel: the bar construction is the
\emph{categorical logarithm} for chiral algebras, mapping
multiplicative structure (the OPE) to additive/nilpotent structure
(the bar differential with $\dzero^2 = 0$).  The cobar construction
is the corresponding \emph{exponential}, reconstructing the algebra
from the nilpotent data.  The structural principle governing the
entire theory is:
\begin{center}
\emph{Nilpotence is the logarithm of periodicity.}
\end{center}
At genus~$0$, the logarithm~$\eta_{12}$ is single-valued
on~$C_2(\mathbb{P}^1)$, the bar differential is nilpotent
($\dzero^2 = 0$), and the monodromy is trivial.  At
genus~$g \geq 1$, the logarithm $d\log E(z_1,z_2)$ acquires
monodromy around the $B$-cycles of~$\Sigma_g$, the fiberwise bar
differential fails nilpotence by
$\dfib^{\,2} = \kappa(\cA)\cdot\omega_g$, and the curvature
$\kappa(\cA)$ is the \emph{infinitesimal generator} of this
monodromy---the residue at the branch point of the logarithm.
This is the chiral Riemann--Hilbert correspondence: the nilpotent
datum $\kappa$ determines the periodic datum (the monodromy of the
bar family over~$\overline{\mathcal{M}}_g$) via exponentiation.

The four theorems are four aspects of this correspondence:
Theorem~A asserts that the categorical logarithm and exponential
exist; Theorem~B that they are mutually inverse on the
convergence domain (the Koszul locus); Theorem~C that the failure
off the convergence domain decomposes into complementary
Lagrangians (the branch structure of the logarithm); and
Theorem~D that the leading coefficient~$\kappa(\cA)$ of the
logarithm determines the full scalar package.  The modular
periodicity profile~$\Pi(\cA)$ of
Remark~\ref{rem:periodicity-triple} is the global/exponential
counterpart of the infinitesimal/nilpotent characteristic
$\kappa(\cA)$: the period~$N$ is the smallest positive integer
for which $\exp(2\pi i N \kappa / 24) = 1$, recovering the
T-matrix periodicity formula of
Conjecture~\ref{thm:modular-periodicity-minimal} as a direct
consequence of the nilpotence-periodicity correspondence.
\end{remark}

Now transplant the construction to genus $g \geq 1$.  On a Riemann
surface $\Sigma_g$, the form $\eta_{12} = d\log(z_1 - z_2)$ is
not well-defined: the function $z_1 - z_2$ has no global meaning.
The replacement is
\begin{equation}\label{eq:genus-g-propagator-intro-ch0}
\eta_{12}^{(g)} = d\log E(z_1, z_2)
+ \sum_{k=1}^{g} \omega_k(z_1) \oint_{A_k} \omega_k(z_2),
\end{equation}
where $E(z_1,z_2)$ is the prime form and $\{\omega_k\}$ is a basis
of $H^0(\Sigma_g, \Omega^1)$, the fiber of the Hodge bundle
$\mathbb{E} \to \overline{\mathcal{M}}_g$ at $[\Sigma_g]$.  The
period correction terms are forced by double-periodicity.  They
couple the collision geometry to the complex structure of the curve,
and this coupling has a consequence: the Arnold relation acquires a
defect,
\begin{equation}\label{eq:curvature-intro}
\dfib^{\,2} = \kappa(\cA) \cdot \omega_g,
\end{equation}
where $\omega_g$ is the Arakelov $(1,1)$-form on $\Sigma_g$ and
$\kappa(\cA)$ is a scalar determined by the leading OPE
singularity.  The bar differential no longer squares to zero.
Period integrals along $A$-cycles restore nilpotence:
$\Dg{g}^{\,2} = 0$ for the total corrected differential
(Theorem~\ref{thm:quantum-diff-squares-zero}).

As $[\Sigma_g]$ varies over the moduli
space~$\overline{\mathcal{M}}_g$, the bar complex traces a family of
curved cochain complexes.  The curvature class
$\mathrm{obs}_g = \kappa(\cA) \cdot \lambda_g \in
H^*(\overline{\mathcal{M}}_g)$, where
$\lambda_g = c_g(\mathbb{E})$ is the top Chern class of the Hodge
bundle, is a characteristic class of this family
(Theorem~\ref{thm:genus-universality}).

The Koszul dual $\cA^!$ traces a dual family, and the two families
are complementary: their obstruction classes pair to give the full
cohomology of moduli with coefficients in the center $Z(\cA)$.
\label{prin:nonabelian-poincare}

\section{Classical Koszul duality}
\label{sec:classical-koszul-confinement}

Classical Koszul duality is a theorem about quadratic algebras on a
point.  Priddy~\cite{Priddy70} constructs, from a quadratic
algebra~$A = T(V)/(R)$, a dual algebra
$A^! = T(V^*)/R^\perp$ and a bar complex
$B(A) = (T^c(s^{-1}\bar{A}), d)$ whose cohomology computes~$A^!$.
When $A$ is Koszul, the cobar inverts the bar:
$\Omega(B(A)) \xrightarrow{\sim} A$.  The theory
(Beilinson--Ginzburg--Soergel~\cite{BGS96},
Loday--Vallette~\cite{LV12}) is complete.

\index{chiral algebra|textbf}
\index{factorization algebra|textbf}
Chiral algebras live on curves.  A chiral algebra on a smooth
algebraic curve~$X$ (Beilinson--Drinfeld~\cite{BD04}) is a
$\mathcal{D}_X$-module whose operations come from the singularity
structure of the OPE\@.  The tensor factors of the bar complex acquire
positions $z_1, \ldots, z_n \in X$, and the configuration space
$C_n(X)$ and its Fulton--MacPherson
compactification~$\overline{C}_n(X)$ enter the construction.

\label{rem:governing-question-route}%
\label{rem:chriss-ginzburg-synthesis}%

\section{The main theorems}
\label{sec:NAP-unifying}

\index{Koszul duality!chiral|textbf}

\begin{theorem}[Unification via configuration spaces; \ClaimStatusProvedHere]
\label{thm:three-way-correspondence}
\index{factorization homology}
\index{Koszul pair!classical}
For a chiral Koszul pair $(\mathcal{A}, \mathcal{A}^!)$ on a
smooth curve~$X$, the three aspects are equivalent:
the bar-cobar adjunction
$\Omega(\bar{B}(\mathcal{A})) \simeq \mathcal{A}$, the
geometric realization of both $\bar{B}(\mathcal{A})$ and
$\mathcal{A}^!$ via integrals on $\overline{C}_n(X)$, and
the factorization homology duality
$\int_X \mathcal{A} \simeq (\int_{-X} \mathcal{A}^!)^\vee$.
The equivalence is the combination of
Theorem~\ref{thm:bar-cobar-isomorphism-main},
Theorem~\ref{thm:higher-genus-inversion}, and
Theorem~\ref{thm:quantum-complementarity-main}.
\end{theorem}
\begin{proof}
The equivalence combines Theorem~A (bar-cobar isomorphism, \S\ref{sec:main-results-complete}), Theorem~B (bar-cobar inversion, ibid.), and Theorem~C (deformation-obstruction complementarity, ibid.): the bar construction provides the geometric bridge (A), inversion provides the algebraic reconstruction (B), and complementarity provides the quantum correction classification (C).
\end{proof}

\begin{remark}[Three-level dictionary: classical $\to$ geometric $\to$ chiral]
\label{rem:three-level-dictionary}
Every chiral Koszul pair arises by enhancement of a classical algebraic
Koszul pair through two functorial stages:
\begin{center}
\begin{tabular}{lll}
\textbf{Level} & \textbf{Setting} & \textbf{Prototype} \\
\hline
Classical & Abstract operads & $\mathrm{Com}^! = \mathrm{Lie}$, $\mathrm{Sym}(V)^! = \Lambda(V^*)$ \\
Geometric & Configuration spaces & Log forms on $\overline{C}_n(\mathbb{R}^d)$ realize bar \\
Chiral & D-modules on curves & $\mathcal{H}_\kappa \leftrightarrow \mathrm{Sym}^{\mathrm{ch}}(V^*)$, $\hat{\mathfrak{g}}_k \leftrightarrow \hat{\mathfrak{g}}^{\!\vee}_{-k-2h^\vee}$ \\
\end{tabular}
\end{center}
Here $\mathfrak{g}^\vee$ denotes the Langlands dual; for simply-laced
$\mathfrak{g}$, one has $\mathfrak{g}^\vee = \mathfrak{g}$.

The passage from classical to geometric is Kontsevich formality (making abstract
structure computable via configuration space integrals); the passage from geometric
to chiral is D-module enhancement (incorporating the factorization axiom and OPE data).
Both stages preserve Koszul duality by
Theorem~\ref{thm:geometric-equals-operadic-bar}.
\end{remark}

\section{The modular Koszul object}
\label{sec:central-thesis}
\label{subsec:structural-perspective}
\index{modular Koszul duality|textbf}
\index{quantum corrections!as modular completion}

A \emph{modular Koszul chiral algebra}
(Definition~\ref{def:modular-koszul-chiral}) is a factorization
algebra $\mathcal{A}$ on $\operatorname{Ran}(X)$ whose duality data
extends over the modular operad of curves, with curvature classes
compatible under clutching, Verdier duality, and trace.  Theorems~A--D
are four aspects of this single structure.

In addition to its Koszul dual~$\mathcal{A}^!$, such an object
carries a \emph{characteristic hierarchy}: a proved scalar package, a
separately proved spectral layer, and a proved full homotopy
completion encoded by a universal Maurer--Cartan class

\begin{equation}\label{eq:universal-MC}
\Theta_{\mathcal{A}} \;\in\;
\operatorname{MC}\!\left(
  \operatorname{Def}_{\mathrm{cyc}}(\mathcal{A})
  \;\widehat{\otimes}\;
  R\Gamma\!\left(\overline{\mathcal{M}}_{g,\bullet},\, \mathbb{Q}\right)
\right)
\end{equation}
(Theorems~\ref{conj:universal-MC} and~\ref{conj:universal-theta}).  The construction proceeds through three inputs --- the intrinsic cyclic $L_\infty$ model (Theorem~\ref{thm:cyclic-linf-graph}), the geometric completed tensor product with modular-operadic clutching (Proposition~\ref{prop:geometric-modular-operadic-mc}), and tautological-line support of the genus-$g$ obstruction (Theorem~\ref{thm:tautological-line-support}) --- assembled in Theorem~\ref{thm:mc2-full-resolution}.

The scalar shadow of this package is the sequence of obstruction
coefficients $\mathrm{obs}_g(\mathcal{A}) = \kappa(\mathcal{A}) \cdot
\lambda_g$ (Theorem~\ref{thm:genus-universality}).  The genus-$1$
coefficient $\kappa(\mathcal{A})$ is not the fundamental concept; it is
the first characteristic number of the universal class
$\Theta_{\mathcal{A}}$ (Theorem~\ref{thm:mc2-full-resolution}).

The four theorems:
\begin{itemize}
\item \emph{Theorem~A} (Geometric bar-cobar duality,
  Theorem~\ref{thm:bar-cobar-isomorphism-main}).
  The reduced bar functor
  \[
  \bar{B}_X \colon \operatorname{Alg}_{\mathrm{aug}}(\operatorname{Fact}(X))
  \to \operatorname{CoAlg}_{\mathrm{conil}}(\operatorname{Fact}(X))
  \]
  is intertwined with Verdier duality:
  $\mathbb{D}_{\operatorname{Ran}} \bar{B}_X(\mathcal{A})
  \simeq \bar{B}_X(\mathcal{A}^!)$,
  functorial in families over $\overline{\mathcal{M}}_{g,n}$.
  Theorem~A decomposes into three layers
  (Remark~\ref{rem:theorem-A-decomposition}):
  the fundamental theorem of chiral twisting morphisms
  ($\mathrm{A}_0$,
  Theorem~\ref{thm:fundamental-twisting-morphisms}),
  bar concentration
  ($\mathrm{A}_1$, Theorem~\ref{thm:bar-concentration}), and
  the Verdier-geometric duality
  ($\mathrm{A}_2$, Theorem~\ref{thm:bar-cobar-isomorphism-main}).

\item \emph{Theorem~B} (Bar-cobar inversion,
  Theorem~\ref{thm:higher-genus-inversion}).
  On the Koszul locus, the counit
  $\Omega_X \bar{B}_X(\mathcal{A}) \to \mathcal{A}$ is an
  equivalence.  Off this locus, the same bar-cobar object persists
  but becomes curved; the failure of inversion is measured by the
  universal Maurer--Cartan class~\eqref{eq:universal-MC}.

\item \emph{Theorem~C} (Deformation-obstruction complementarity,
  Theorem~\ref{thm:quantum-complementarity-main}).
  The ambient complex
  $\mathbf{C}_g(\cA) := R\Gamma(\overline{\mathcal{M}}_g, \mathcal{Z}_{\cA})$
  carries a shifted-symplectic pairing from Verdier duality on the
  center local system.  The Verdier involution~$\sigma$ decomposes
  $\mathbf{C}_g(\cA)$ into homotopy eigenspaces
  $\mathbf{Q}_g(\cA) \oplus \mathbf{Q}_g(\cA^!)$
  (Definition~\ref{def:complementarity-complexes}),
  which embed as complementary Lagrangians.
  At the cohomological level,
  $Q_g(\cA) \oplus Q_g(\cA^!) \cong
  H^*(\overline{\mathcal{M}}_g, \mathcal{Z}_{\cA})$.
  Theorem~C decomposes into two layers
  (Remark~\ref{rem:theorem-C-decomposition}):
  the fiber--center identification
  ($\mathrm{C}_0$,
  Theorem~\ref{thm:fiber-center-identification}), which
  produces $\mathcal{Z}_{\cA}$ from the relative bar family,
  and the chain-level Lagrangian polarization
  ($\mathrm{C}_1$,
  Theorem~\ref{thm:quantum-complementarity-main}).

\item \emph{Theorem~D$_{\mathrm{scal}}$} (Scalar modular characteristic,
  Theorem~\ref{thm:modular-characteristic}).
  A single scalar $\kappa(\cA)$---the genus-$1$ curvature
  coefficient---determines the entire scalar modular package
  (Definition~\ref{def:scalar-modular-package}): $\kappa$ is
  universal ($\mathrm{obs}_g = \kappa \cdot \lambda_g$),
  additive under tensor product, and anti-symmetric under Koszul
  duality ($\kappa + \kappa' = 0$ for affine Kac--Moody; more
  generally $\kappa + \kappa' = \varrho(\fg) \cdot K$ for
  W-algebras).  Its vanishing is anomaly
  cancellation; its generating function is the $\hat{A}$-genus.

  \emph{Theorem~D$_\Delta$} (Spectral characteristic,
  Theorem~\ref{thm:spectral-characteristic}):
  the spectral discriminant $\Delta_{\cA}(x)$ is a
  separately proved non-scalar invariant, depending on the
  full quadratic OPE data.

  The universal Maurer--Cartan class $\Theta_{\cA}$ is proved
  (Theorem~\ref{conj:universal-theta},
  Definition~\ref{def:full-modular-package}) and is
  scalar-saturated: $\Theta_{\cA}^{\min} = \kappa(\cA) \cdot
  \eta \otimes \Lambda$
  (Corollary~\ref{cor:scalar-saturation}).
  The saturation equivalence criterion
  (Proposition~\ref{prop:saturation-equivalence}) and
  functorial stability
  (Proposition~\ref{prop:saturation-functorial}) reduce the
  characteristic quintuple to an effective quadruple for
  the entire standard Lie-theoretic landscape
  (Corollary~\ref{cor:effective-quadruple}).
  At generic level, cyclic rigidity
  (Theorem~\ref{thm:cyclic-rigidity-generic}) proves the
  quadruple reduction for \emph{all} vertex algebras with
  simple symmetry via Whitehead's lemma and
  Kazhdan--Lusztig semisimplicity.
  Conjecture~\ref{conj:scalar-saturation-universality}
  extends this to admissible levels.
\end{itemize}

\begin{remark}[Physical meaning of the four theorems]
\label{rem:trilogy}
\index{trilogy reading}

\smallskip\noindent\emph{Theorem~A} (bar-cobar duality).
\textbf{Mathematics}: the Verdier involution
$\mathbb{D}_{\operatorname{Ran}}$ intertwines the bar functor with
the Koszul dual's bar functor---an equivalence of functors on
augmented factorization algebras.
\textbf{Mathematical physics}: the bar complex is the
BRST complex of the gauge-fixed chiral algebra, and Verdier
duality is the CPT involution on the field content
(Theorem~\ref{thm:qme-bar-cobar}).
\textbf{Physics}: the Koszul dual is the ghost system of the
matter algebra; bar-cobar duality is the statement that
gauge-fixing is involutive.

\smallskip\noindent\emph{Theorem~B} (inversion).
\textbf{Mathematics}: the counit
$\Omega_X \bar{B}_X(\cA) \xrightarrow{\sim} \cA$ is a
quasi-isomorphism on the Koszul locus; the spectral sequence
collapses at~$E_2$.
\textbf{Mathematical physics}: the path integral over the
BRST complex reconstructs the original theory; the $E_2$
collapse is the one-loop exactness of the ghost OPE\@.
\textbf{Physics}: the BRST resolution is complete---every
physical state is recovered from bar-complex data.

\smallskip\noindent\emph{Theorem~C} (complementarity).
\textbf{Mathematics}: the ambient complex $\mathbf{C}_g(\cA)$
carries a Lagrangian decomposition into deformation and
obstruction eigenspaces under the Verdier involution.
\textbf{Mathematical physics}: what one algebra sees as a
deformation of the modular structure, its Koszul dual sees as
an obstruction; the Kodaira--Spencer map connects the two.
\textbf{Physics}: anomaly cancellation $\kappa = 0$ is
mutual consistency of matter and ghost; the anomaly of one
is the ghost-number violation of the other.

\smallskip\noindent\emph{Theorem~D$_{\mathrm{scal}}$}
(scalar characteristic).
\textbf{Mathematics}: a single scalar $\kappa(\cA)$
determines genus expansions, the family index, and the
$\hat{A}$-genus generating function.
\textbf{Mathematical physics}: the genus-$g$ curvature
$\dfib^{\,2} = \kappa \cdot \omega_g$ is a Chern class on
$\overline{\mathcal{M}}_g$ controlled by one universal number.
\textbf{Physics}: $\kappa$ is the one-loop coefficient of the
effective action; the $\hat{A}$-genus is the gravitational
anomaly polynomial; $\kappa = 0$ is anomaly-free string
propagation at all genera.
\end{remark}

\begin{remark}[Construction versus resolution: the convergence radius of the logarithm]
\label{rem:construction-vs-resolution-intro}
\label{rem:constitutional-reading-rule}
\index{construction versus resolution}
For \emph{every} augmented chiral algebra~$\cA$, the bar complex
$\bar{B}_X(\cA)$ and the cobar construction
$\Omega_X(\bar{B}_X(\cA))$ exist as objects
(Definition~\ref{def:chiral-twisting-datum}).  But the counit
$\Omega_X(\bar{B}_X(\cA)) \to \cA$ is a quasi-isomorphism only when
the canonical twisting morphism is Koszul
(Definition~\ref{def:chiral-koszul-morphism}).  Thus Theorem~A
(duality) holds for all augmented chiral algebras, while Theorem~B
(inversion) holds on the Koszul locus.

In the language of Remark~\ref{rem:nilpotence-periodicity}, the
Koszul locus is the \emph{convergence radius of the categorical
logarithm}: the bar complex ``converges'' (the PBW spectral sequence
degenerates at~$E_2$) precisely where the logarithm faithfully
encodes the algebra.  Off this locus, the logarithm ``diverges'' in
the ordinary derived category; the coderived category of
Positselski~\cite{Positselski11} provides the analytic continuation
that extends the bar-cobar correspondence past the convergence radius.
At higher genus with $\kappa \neq 0$, the curved fiberwise differential
$\dfib^{\,2} = \kappa \cdot \omega_g$ forces the logarithm into its
multi-valued regime, and the coderived framework becomes essential:
the underlying cochain complex of the curved bar complex is acyclic,
but the coderived category retains the full homotopical content---just
as the logarithm function is trivial on a simply-connected domain but
carries non-trivial monodromy data on a multiply-connected one.

The characteristic invariants---scalar ($\kappa$), spectral
($\Delta_\cA$), and Maurer--Cartan class~$\Theta_\cA$---form a
hierarchy of increasing refinement
(Remark~\ref{rem:characteristic-hierarchy}), measuring the
logarithm at successively finer levels.
\end{remark}

\noindent
The genus tower $\{\bar{B}^{(g)}(\mathcal{A})\}_{g \geq 0}$ is a
Maurer--Cartan deformation of the genus-$0$ bar differential inside a
cyclic $L_\infty$ algebra
(Definition~\ref{def:cyclic-deformation-elementary}).  Its most
concrete shadow is the obstruction class
$\mathrm{obs}_g = \kappa(\cA) \cdot \lambda_g$ in the tautological
ring of~$\overline{\mathcal{M}}_g$.  At chain level, the deformation
is a curved coderivation $m_0^{(g)}$ in the bar-cobar package.  The
conjectural completion is a universal Maurer--Cartan element
$\Theta_{\cA} \in
\operatorname{MC}(\operatorname{Def}_{\mathrm{cyc}}(\cA)
\;\widehat{\otimes}\;
R\Gamma(\overline{\mathcal{M}}_{g,\bullet}, \mathbb{Q}))$
from which the chain-level and scalar data are recovered
(Theorem~\ref{conj:universal-theta}).

The scalar and chain-level packages are proved.  The proved families
include the free fields
(Proposition~\ref{prop:standard-examples-modular-koszul}), affine
Kac--Moody algebras at all genera
(Theorem~\ref{thm:pbw-allgenera-km},
Corollary~\ref{cor:unconditional-allgenera-km}), the Virasoro family at
all genera
(Theorem~\ref{thm:pbw-allgenera-virasoro},
Corollary~\ref{cor:unconditional-allgenera-virasoro}), and principal
finite-type $\mathcal{W}$-families, including higher $\mathcal{W}_N$
(Theorem~\ref{thm:pbw-allgenera-principal-w},
Corollary~\ref{cor:unconditional-allgenera-principal-w}).

The proved scalar and spectral characteristic hierarchy already
explains, in one move, the shared
discriminant phenomenon
(Theorem~\ref{thm:ds-bar-gf-discriminant}), the complementarity sums
$c + c' = K$ (Theorem~\ref{thm:central-charge-complementarity}), the
universal genus generating function
(Theorem~\ref{thm:universal-generating-function}), and the
$\hat{A}$-genus appearance in genus expansions.

Among the shadows of~$\Theta_{\mathcal{A}}$, the \emph{spectral
discriminant}
\[
\Delta_{\mathcal{A}}(x)
\;:=\; \operatorname{disc}\bigl(P_{\mathcal{A}^!}(x)\bigr),
\]
the discriminant of the Koszul dual Hilbert series as an algebraic
function, deserves special mention.  It is the first genuinely
non-scalar characteristic invariant of the modular Koszul object:
invariant under Koszul duality ($\Delta_{\mathcal{A}} =
\Delta_{\mathcal{A}^!}$, Proposition~\ref{prop:discriminant-characteristic})
and under Drinfeld--Sokolov reduction
(Theorem~\ref{thm:ds-bar-gf-discriminant}).  The shared discriminant
$\Delta(x) = (1-3x)(1+x)$ across the $\widehat{\mathfrak{sl}}_2$,
Virasoro, and $\beta\gamma$ families is not a coincidence but a
consequence: these algebras lie on the same modular spectral sheet,
even when their local operator content is entirely different.

The full characteristic package of a modular Koszul chiral algebra is
\[
\mathcal{C}_{\cA}
=
(\Theta_{\cA}, \kappa(\cA), \Delta_{\cA}, \Pi_{\cA}, \mathcal{H}_{\cA}),
\]
a hierarchy of increasing refinement: the scalar $\kappa(\cA)$,
the spectral discriminant $\Delta_{\cA}(x)$, and the
universal class $\Theta_{\cA}$
\textup{(}Definition~\ref{def:full-modular-package}\textup{)}.

The universal class $\Theta_\cA$ is constructed by
Theorem~\ref{thm:mc2-full-resolution}, resolving the original
Conjecture~\ref{conj:master-theta} (MC2) through three inputs:
the intrinsic cyclic model, the geometric completed
tensor product with clutching, and a one-channel normalization in the
simple-Lie case
\textup{(}Proposition~\ref{prop:mc2-reduction-principle},
Corollary~\ref{cor:one-dim-obstruction};
see \S\ref{sec:modular-koszul-programme} for the full reduction\textup{)}.
The factorization-categorical
(Conjecture~\ref{conj:master-dk-kl}) and BV/BRST
(Conjecture~\ref{conj:master-bv-brst}) identifications are downstream.

\begin{remark}[Five irreducible ingredients]\label{rem:four-pieces}
\index{modular Koszul duality!five irreducible pieces}
\index{double-frame architecture|textbf}
Modular Koszul duality rests on five concrete geometric ingredients.
The first four belong to the commutative ($\Einf$) face; the fifth
belongs to the associative ($\Eone$) face.
\begin{enumerate}[label=\textup{(\roman*)}]
\item \emph{Three-point collision.}
  The Arnold relation $\omega_{12} \wedge \omega_{23}
  + \omega_{23} \wedge \omega_{31} + \omega_{31} \wedge \omega_{12} = 0$
  (Theorem~\ref{thm:arnold-relations}) is the coherence statement
  that fusing three insertions does not depend on the order of pairwise
  collisions.  In the $\operatorname{Ran}$-space language, this is
  factorization coherence itself.  It is the genus-$0$ seed from which
  all later duality grows.
\item \emph{Verdier duality on $\operatorname{Ran}(X)$.}
  The bar-cobar adjunction is controlled by Verdier duality:
  \[
  \mathbb{D}_{\operatorname{Ran}} \bar{B}_X(\cA)
  \simeq \bar{B}_X(\cA^!)
  \]
  (Theorem~\ref{thm:bar-cobar-isomorphism-main}).
  This makes the Koszul dual a theorem rather than a
  definition-by-analogy; at chain level, it is the shadow of
  non-abelian Poincar\'e duality in the sense of
  Ayala--Francis.
\item \emph{Genus-$1$ curvature.}
  Once the propagator acquires periods, the differential develops
  fiberwise curvature $\dfib^{\,2} = \kappa(\cA) \cdot \omega_1 \cdot \operatorname{id}$
  (Theorem~\ref{thm:genus-universality};
  Convention~\ref{conv:higher-genus-differentials}).  The natural ambient world
  is Positselski's coderived category, not the ordinary derived
  category: quantum corrections are a change in the ambient
  homological category, not extra terms in a formula.
\item \emph{Clutching of stable curves.}
  The genus tower $\{\bar{B}^{(g)}(\cA)\}_{g \geq 0}$ is
  compatible with boundary clutching maps via the modular operad
  structure (Theorem~\ref{thm:prism-higher-genus}).  This is where
  the theory becomes genuinely modular: the relevant combinatorics
  is stable graphs, not just trees, and the conjectural universal class
  $\Theta_{\cA}$ (Theorem~\ref{conj:universal-MC})
  would be compatible with clutching, trace, and
  Verdier duality.
\item \emph{Associative structure and ordered factorization.}
  \index{ordered factorization|see{$\Eone$-chiral algebra}}
  The Yangian $Y(\mathfrak{g})$ is $\Eone$-chiral
  (Definition~\ref{def:e1-chiral}): associative but not commutative,
  so its bar complex uses \emph{ordered} configuration spaces.
  Verdier duality acts by $R$-matrix inversion $R \mapsto R^{-1}$
  on the module category
  (Theorem~\ref{thm:yangian-koszul-dual}), not by level shift.
  The sign flip $\hbar \mapsto -\hbar$ replaces the
  Feigin--Frenkel involution of the commutative face.  This piece is
  logically independent of genus-$1$ curvature: it appears already at
  genus~$0$ and forces the theory to accommodate both commutative and
  associative chiral algebras within a single framework.
\end{enumerate}
The five pieces organize into a double-frame architecture.

\smallskip\noindent\textbf{Frame~A: commutative/modular}
(pieces (i)--(iv), exhibited by the Heisenberg algebra in
Chapter~\ref{ch:heisenberg-frame}).
One current, one double pole, yet all four mechanisms visible: Arnold
ensures $d^2 = 0$; Verdier produces the dual
$\mathcal{H}_k^! = \mathrm{Sym}^{\mathrm{ch}}(V^*)$; the genus-$1$
period integral creates curvature $\kappa = 1/k$; clutching assembles
the genus tower.  The Heisenberg algebra makes modular completion
unavoidable.

\smallskip\noindent\textbf{Frame~B: associative/factorization}
(pieces (i), (ii), (v), exhibited on the Yangian evaluation locus in
Chapter~\ref{chap:yangians}).
The same geometric engine---logarithmic forms on Fulton--MacPherson
compactifications, Verdier duality, Borcherds extraction---operates
on \emph{ordered} configurations, producing $R$-matrix inversion on
the module category and the derived Drinfeld--Kohno bridge
(Theorem~\ref{thm:factorization-dk-eval}).  The Yangian evaluation
locus makes $\Eone$-factorization unavoidable; the $R$-matrix
structure on its representations makes quantum-group duality visible.

\smallskip\noindent
The general theory is not an arbitrary abstraction imposed from
above.  It is \emph{forced} by the demand that these two faces
coexist: a theory of Koszul duality on curves must simultaneously
handle the commutative regime (where genus curvature and clutching
produce the modular package) and the associative regime (where
noncommutativity forces ordered configurations and Verdier duality
produces $R$-matrix inversion on the module category).  That these
two structurally distinct demands are resolved by the same geometric
mechanism---Verdier duality on configuration space
compactifications, ordered or unordered---is the central structural
insight of this monograph.
\end{remark}

\begin{remark}[The DK square as witness of inevitability]
\label{rem:dk-inevitability}
\index{Drinfeld--Kohno!derived!as double-frame witness}
The double-frame architecture is not a pedagogical device.
It is witnessed mathematically by the derived Drinfeld--Kohno
square (Theorem~\ref{thm:derived-dk-affine}):
\begin{equation}\label{eq:dk-intro}
\begin{tikzcd}
\mathcal{O}_k^{\mathrm{int}}(\widehat{\mathfrak{g}})
  \ar[r, "\Phi", "\sim"']
  \ar[d, "\Phi_{\mathrm{KL}}"', "\sim"]
& \mathcal{O}_{k'}(\widehat{\mathfrak{g}})
  \ar[d, "\Phi_{\mathrm{KL}}", "\sim"']
\\
\mathcal{C}(U_q(\mathfrak{g}))
  \ar[r, "q \mapsto q^{-1}"', "\sim"]
& \mathcal{C}(U_{q^{-1}}(\mathfrak{g}))
\end{tikzcd}
\end{equation}
The top arrow is bar-cobar duality on affine Kac--Moody module
categories (Frame~A: level inversion $k \mapsto k' = -k - 2h^\vee$,
the $\Einf$-chiral Verdier exchange).
The bottom arrow is $R$-matrix inversion
$R_q \mapsto R_q^{-1}$ (Frame~B: the quantum-group involution
$q \mapsto q^{-1}$).  The vertical arrows are the Kazhdan--Lusztig
equivalences, which identify conformal-block monodromy with the
quasi-$R$-matrix.  The commutativity of the square is the theorem
that bar-cobar duality \emph{geometrically implements} the quantum
group involution: orientation reversal on configuration spaces
realizes $q \mapsto q^{-1}$.

This is the precise sense in which the two frames cannot be studied
independently.  Any theory of Koszul duality on curves that handles
the commutative face (level inversion, modular curvature) must also
handle the associative face ($R$-matrix inversion on module
categories), because they are connected by the KL equivalences.
The DK square is the mathematical certificate that a unified
framework---modular Koszul duality---is forced, not chosen.

The scalar modular characteristic (Theorem~D$_{\mathrm{scal}}$)
is the trace of this square.  The top arrow sends
$k \mapsto k' = -k - 2h^\vee$, so
$\kappa(\cA^!) = \kappa(\widehat{\fg}_{k'})
= \frac{(k' + h^\vee)\dim\fg}{2h^\vee}
= -\kappa(\cA)$:
the anti-symmetry $\kappa + \kappa' = 0$ is the scalar shadow
of the DK square's top arrow.  On the quantum-group side, the
same operation is $q \mapsto q^{-1}$.  The four theorems and
the DK square are therefore not separate achievements;
they are four projections of one geometric structure on
$\operatorname{Ran}(X) \times \overline{\mathcal{M}}_g$.
\end{remark}

\begin{remark}[Scope]\label{rem:two-strata}
Theorems~A--D, the Lagrangian refinement of complementarity, the
scalar and spectral characteristic packages, and unconditional modular
Koszulity for free fields, Kac--Moody, Virasoro, and principal
$\mathcal{W}$-algebras are proved.  The universal Maurer--Cartan
class $\Theta_{\cA}$ (Theorem~\ref{conj:universal-MC}), the full
coderived factorization formalism, and the categorical Drinfeld--Kohno
extension beyond the evaluation locus remain open.  The precise
frontier is recorded in Chapter~\ref{chap:concordance} and the
programme of~\S\ref{sec:modular-koszul-programme}.
\end{remark}

\begin{convention}[Semantic levels: homotopy, model, shadow]
\label{conv:hms-levels}
\index{semantic levels|textbf}
\index{H/M/S convention|see{semantic levels}}
Three levels of description recur throughout.

\smallskip\noindent\textbf{H-level (homotopy/native).}
The statement in the stable $\infty$-category of factorization
(co)algebras on $\operatorname{Ran}(X)$, or in the associated
formal moduli problem / shifted-symplectic language.  At this
level the bar construction is an object
$\mathcal{B}_X(\cA) \in
\operatorname{CoAlg}_{\mathrm{conil}}(\operatorname{Fact}(X))$,
defined up to contractible ambiguity.

\smallskip\noindent\textbf{M-level (model).}
An explicit dg / $L_\infty$ / filtered complex presentation.
The geometric bar complex
$\bar{B}_X(\cA) = (T^c(s^{-1}\bar{\cA}), d)$ built from
configuration space integrals on $\overline{C}_n(X)$ is the
canonical model throughout this monograph.

\smallskip\noindent\textbf{S-level (shadow).}
The cohomological, Ext, numerical, or generating-function
consequence: bar cohomology dimensions, Hilbert series,
obstruction scalars $\kappa(\cA)$, free energies $F_g(\cA)$.

\smallskip\noindent\textbf{Notation rules.}
\begin{itemize}
\item $\simeq$\,: equivalence in the homotopy/$\infty$-categorical
  sense (H-level).
\item $\cong$\,: strict isomorphism of chosen dg models (M-level).
\item $=$\,: definition or literal equality only.
\end{itemize}
The proved results (Remark~\ref{rem:two-strata}) are proved at
M-level and stated at H-level.

\smallskip\noindent\textbf{Model independence.}
Any two admissible dg presentations of the same factorization object
are connected by a contractible space of quasi-isomorphisms
(Proposition~\ref{prop:model-independence}).  This lemma is the
constructive content of ``beyond chain level'': it ensures that
M-level computations are legitimate shadows of H-level theorems,
without recourse to full
$\infty$-categorical abstraction.
\end{convention}

\label{conv:proof-architecture}%

\begin{convention}[Regime tags]\label{conv:regime-tags}
\index{regime tag}
Theorems in this monograph operate within one of four \emph{regimes},
listed in order of increasing generality:
\begin{enumerate}
\item \textbf{Quadratic.}  The chiral algebra $\cA$ is generated by fields
  whose OPE relations are quadratic (simple poles only).  This is the
  classical Koszul setting: the bar differential has no curvature
  ($m_0 = 0$, $\dzero^2 = 0$), and Theorems~A--D hold without
  qualification.  Examples: Heisenberg, free fermions, lattice VOAs
  at self-dual radius.
\item \textbf{Curved-central.}  The genus-$g$ bar complex acquires
  fiberwise curvature $\dfib^{\,2} = \kappa \cdot \omega_g$ controlled
  by a single scalar $\kappa \in k$ (the scalar modular
  characteristic).
  The total differential $\Dg{g}$ is still square-zero.  This is the
  generic interacting setting for Koszul algebras.
  Examples: $\widehat{\mathfrak{g}}_k$ at non-critical level, Virasoro
  with $c \neq 0$.
\item \textbf{Filtered-complete.}  The algebra carries a complete
  decreasing filtration $F^\bullet \cA$ with $\operatorname{gr}^F \cA$
  quadratic Koszul.  Bar-cobar inversion holds in the filtered-complete
  sense (Theorem~\ref{thm:higher-genus-inversion}).
  Examples: $W$-algebras via DS reduction, deformation quantizations.
\item \textbf{Programmatic.}  Results at this level are part of the
  modular homotopy programme (Stratum~II,
  Remark~\ref{rem:two-strata}) and are stated as conjectures with
  identified gaps.  Examples: derived Drinfeld--Kohno,
  $\mathbb{E}_n$ Koszul duality for $n \geq 2$.
\end{enumerate}
When a theorem's hypotheses place it in a specific regime, the regime
tag appears in the statement.  Results without an explicit tag hold
in all regimes where the hypotheses are satisfied.

\smallskip\noindent
\textbf{Regime hierarchy as convergence of the categorical logarithm.}
In the language of Remark~\ref{rem:nilpotence-periodicity}, the four
regimes correspond to increasing levels of non-trivial monodromy of the
categorical logarithm (the bar construction).  In the quadratic regime,
the logarithm is single-valued: $\dzero^2 = 0$ and the monodromy is
trivial.  In the curved-central regime, the logarithm acquires
monodromy controlled by a single scalar~$\kappa$: $\dfib^{\,2} =
\kappa \cdot \omega_g$, but the corrected logarithm $\Dg{g}$ is
single-valued on the universal cover.  In the filtered-complete regime,
the logarithm converges only after filtration-completion: the PBW
spectral sequence requires multiple pages, and the convergence domain
(Koszul locus) is reached in the associated graded.  In the
programmatic regime, the full monodromy structure (including
categorical and coderived data) is not yet established.
\end{convention}

\medskip

% [Relationship to foundational work (Beilinson--Drinfeld, Costello--Gwilliam,
%  Ayala--Francis): see the concordance chapter for a detailed comparison.]

\section{The total space}
\label{sec:modular-homotopy-theory}
\index{modular homotopy theory|textbf}
\index{total space|textbf}

The nilpotence-periodicity correspondence of
Remark~\ref{rem:nilpotence-periodicity} becomes geometric in the total
space of the universal configuration fibration over moduli.
Consider the universal configuration space over moduli:
\begin{equation}\label{eq:total-fibration-intro}
\pi \colon \overline{C}_n(\mathcal{C}_g) \;\longrightarrow\;
\overline{\mathcal{M}}_g.
\end{equation}
Its fiber over $[\Sigma_g] \in \overline{\mathcal{M}}_g$ is
$\overline{C}_n(\Sigma_g)$, the compactified configuration space of
$n$ distinct points on $\Sigma_g$.  The boundary of this total space
has two types of divisors:
\begin{enumerate}[label=\textup{(\alph*)}]
\item \emph{Collision divisors} $D_I$ (vertical): a subset of marked
  points collides on a fixed smooth curve.  These produce the bar
  differential via residues.
\item \emph{Degeneration divisors} $\delta_\Gamma$ (horizontal): the
  underlying curve acquires a node.  These produce the period
  corrections that restore $\Dg{g}^{\,2} = 0$.
\end{enumerate}
At genus~$0$, only type~(a) exists and the Arnold relation alone
gives~$\dzero^2 = 0$.  At genus~$g \geq 1$, both types coexist.
The Leray spectral sequence of~$\pi$ and the genus spectral sequence
from $F^k = \bigoplus_{g \geq k} \barB^{(g)}(\cA)$ decompose
the total differential $\Dg{g}$ along these two boundary types.

\subsection{The Hodge bundle and the propagator}
\label{subsec:hodge-deformation-engine}

The Hodge bundle $\mathbb{E} = R^1\pi_*\mathcal{O}_{\mathcal{C}_g}$
over $\overline{\mathcal{M}}_g$ has fiber
$H^0(\Sigma_g, \Omega^1)$ over $[\Sigma_g]$: the space of holomorphic
differentials.  The genus-$g$ propagator~\eqref{eq:genus-g-propagator-intro-ch0}
depends on a basis $\{\omega_k\}_{k=1}^g$ of this fiber---a frame for
$\mathbb{E}^\vee|_{[\Sigma_g]}$.  The period corrections
$\sum_k \omega_k(p_i) \oint_{A_k} \omega_k(p_j)$ live in
$\mathbb{E}^\vee \boxtimes \mathbb{E}$.

The Arnold defect~\eqref{eq:curvature-intro} is
$\dfib^{\,2} = \kappa(\cA) \cdot \omega_g$, where $\omega_g$ is the
Arakelov form.  The fiberwise curvature coderivation defines a section
of $\operatorname{End}(\barB(\cA)) \otimes \mathbb{E}^\vee$.  Taking
the scalar trace over the top exterior power:
\begin{equation}\label{eq:lambda-from-trace}
\mathrm{obs}_g(\cA)
= \mathrm{tr}\bigl(\kappa(\cA) \cdot \omega_g\bigr)
= \kappa(\cA) \cdot c_g(\mathbb{E})
= \kappa(\cA) \cdot \lambda_g
\;\in\; H^*(\overline{\mathcal{M}}_g).
\end{equation}
The class $\lambda_g$ is the top Chern class of the Hodge bundle:
the determinant of the space of propagator corrections.

\subsection{The family index}
\label{subsec:family-index-intro}
\index{A-hat genus@$\hat{A}$-genus!as family index}

The family $\barB_g(\cA) \to \overline{\mathcal{M}}_g$ of bar
complexes defines a virtual vector bundle
$[\barB_g(\cA)]^{\mathrm{vir}} \in K^0(\overline{\mathcal{M}}_g)$.
Apply the Grothendieck--Riemann--Roch theorem to the universal
fibration $\pi \colon \overline{\mathcal{C}}_g \to
\overline{\mathcal{M}}_g$:
\[
\mathrm{ch}\bigl([\barB_g(\cA)]^{\mathrm{vir}}\bigr)
= \pi_*\!\bigl(\mathrm{ch}(\barB(\cA)|_{\mathrm{fiber}})
\cdot \mathrm{Td}(T_\pi)\bigr).
\]
The fiber Chern character contributes $\kappa(\cA)$; the Todd class
of the universal curve produces the $\hat{A}$-series.  The
generating function of obstruction classes is
$\sum_{g \geq 1} \mathrm{obs}_g \cdot \hbar^{2g-2}
= \kappa(\cA) \cdot (\hat{A}(i\hbar) - 1)$
(Theorem~\ref{thm:family-index}).  For a free fermion, this is the
spin partition function; for the Virasoro algebra, the one-loop
string amplitude; for $\widehat{\fg}_k$, the leading Verlinde
coefficient.

\subsection{The Lagrangian decomposition}
\label{subsec:verdier-symplectic}
\index{Verdier duality!as symplectic structure}

Verdier duality $\sigma = \mathbb{D}_{\operatorname{Ran}}$ on
$\operatorname{Ran}(X)$ acts on the global sections
$\mathbf{C}_g(\cA) = R\Gamma(\overline{\mathcal{M}}_g,
\mathcal{Z}_{\cA})$ of the center local system.  The composition
\[
\mathbf{C}_g(\cA) \otimes \mathbf{C}_g(\cA)
\;\xrightarrow{\;\mathrm{id} \otimes \sigma\;}
\mathbf{C}_g(\cA) \otimes \mathbf{C}_g(\cA^!)
\;\xrightarrow{\;\mathrm{eval}\;}
R\Gamma(\overline{\mathcal{M}}_g, \omega_{\overline{\mathcal{M}}_g})
\]
is a $(-(3g{-}3))$-shifted symplectic form
(Pantev--To\"en--Vaqui\'e--Vezzosi~\cite{PTVV13}).  The involution
$\sigma$ decomposes $\mathbf{C}_g(\cA)$ into $\pm 1$ eigenspaces
$\mathbf{Q}_g(\cA) \oplus \mathbf{Q}_g(\cA^!)$.  The symplectic
form pairs these eigenspaces perfectly: each is Lagrangian.  Verdier
duality on $\overline{C}_n(\Sigma_g)$ exchanges residues (bar) with
delta functions (cobar); traced through~$\pi$ and projected to
moduli, this exchange produces the complementary Lagrangian
decomposition of Theorem~C.

\section{Dictionary: chiral algebra types and three Koszul dualities}
\label{sec:dictionary}

Collision geometry on a curve forces three levels of algebraic
structure --- commutative ($\Einf$), associative ($\Eone$),
semiclassical ($\Pinf$) --- each carrying its own Koszul duality.

\begin{definition}[$\Einf$-chiral algebra = vertex algebra]\label{def:einf-chiral}
\index{$\mathbb{E}_\infty$-chiral algebra|textbf}
Throughout this monograph, we work over $k = \mathbb{C}$ and let $X$
denote a smooth algebraic curve over~$\mathbb{C}$.  The ambient
$\infty$-category is $\mathcal{V} = \mathcal{D}\text{-mod}(X)$, the
stable $\infty$-category of (right) $\mathcal{D}_X$-modules, equipped
with the \emph{chiral tensor product}
$\otimes^{\mathrm{ch}}$ \cite[\S3.4]{BD04}.  Duality in this category
means \emph{Verdier duality} $\mathbb{D}$ of holonomic
$\mathcal{D}$-modules (not the na\"{\i}ve linear dual
$\mathrm{Hom}_k(-, k)$); see
Remark~\ref{rem:verdier-vs-linear} below.

An \emph{$\Einf$-chiral algebra} on $X$ is a commutative
algebra object in the chiral tensor category
$(\mathcal{D}\text{-mod}(X),\allowbreak
\otimes^{\mathrm{ch}})$
(cf.~\cite[§5.1]{HA} for $E_\infty$-algebras
in symmetric monoidal $\infty$-categories).
Equivalent formulations:
\begin{enumerate}
\item[(a)] A $\mathcal{D}_X$-module with chiral Lie bracket
$\mu\colon j_*j^*(\mathcal{A} \boxtimes \mathcal{A})
\to \Delta_!\mathcal{A}$
\cite[\S3.3]{BD04}, equivalent to the commutative viewpoint via the
chiral envelope \cite[Theorem~3.4.9]{BD04}.
\item[(b)] A vertex algebra $(V,Y,\mathbf{1},T)$ with the locality axiom
\cite[\S3.7]{FBZ04} (requires $X = \mathbb{A}^1$).
\item[(c)] A factorization algebra on $\operatorname{Ran}(X)$
(Definition~\ref{def:factorization-algebra-AF}).
\end{enumerate}
The subscript $\infty$ refers to the factorization level,
not the little discs operad \cite[§5.1]{HA}.
\end{definition}

\begin{remark}[Verdier duality versus linear duality]
\label{rem:verdier-vs-linear}
\index{Verdier duality}
Throughout this monograph, for a holonomic $\mathcal{D}_X$-module
$\mathcal{M}$, we write $\mathbb{D}(\mathcal{M})$ for the
\emph{Verdier dual} $\mathcal{D}$-module (the derived dual in the
category of $\mathcal{D}$-modules, with the appropriate shift by
$\dim X$).  This is \emph{not} the na\"{\i}ve linear dual
$\mathcal{M}^\vee := \mathrm{Hom}_k(\mathcal{M}, k)$, which is
ill-behaved for infinite-dimensional modules.  When we write
$\mathcal{A}^!$ for the Koszul dual, the underlying duality used
in its construction is always Verdier duality on the configuration
spaces $\overline{C}_n(X)$ (see
Theorem~\ref{thm:bar-cobar-verdier}).
\end{remark}

\begin{definition}[$\Eone$-chiral algebra = nonlocal vertex algebra]\label{def:e1-chiral}
\index{$\mathbb{E}_1$-chiral algebra|textbf}
An \emph{$\Eone$-chiral algebra} on $X$ is an associative (but not necessarily commutative)
algebra object in the chiral compound tensor $\infty$-category of $\mathcal{D}_X$-modules.
Concretely, it is an algebra over the \emph{associative chiral operad} $\chirAss$:
a $\mathcal{D}_X$-module $\mathcal{A}$ with a binary chiral operation
\[
\mu: \mathcal{A} \chirtensor \mathcal{A} \to \Delta_!\mathcal{A}
\]
satisfying $\mu \circ (\mu \otimes \operatorname{id}) = \mu \circ (\operatorname{id} \otimes \mu)$
but \emph{not} necessarily skew-symmetry.

In vertex algebra language (for $X = \mathbb{A}^1$ with translation), this
corresponds to a \emph{nonlocal vertex algebra} in the sense of Li~\cite{Li96}
(cf.\ \cite[§4.1]{HA} for the general theory of $\Eone$-algebras):
a quadruple $(V, Y, \mathbf{1}, T)$ satisfying the vacuum axiom, translation
equivariance, and weak associativity, but \emph{not} the locality
(mutual commutativity) axiom.  The precise equivalence is due to
Beilinson--Drinfeld \cite[Theorem~3.4.9]{BD04}
and Francis--Gaitsgory \cite{FG12}.
\end{definition}

\begin{definition}[The associative chiral operad]\label{def:chiral-ass-operad}
The \emph{associative chiral operad} $\chirAss$ on a smooth curve $X$ is the
non-$\Sigma$ chiral operad with $\chirAss(n) = \Delta^{(n)}_!\, \omega_X$
(supported on the small diagonal, without $\Sigma_n$-action).
The quadratic presentation is
$\chirAss = \operatorname{Free}^{\mathrm{ch}}(\mu) / (\mu \circ_1 \mu = \mu \circ_2 \mu)$,
the chiral lift of $\operatorname{Ass}$.
\end{definition}

\begin{proposition}[$\chirAss$ self-duality; \ClaimStatusProvedHere]
\label{prop:chirAss-self-dual}
The associative chiral operad is Koszul self-dual:
\[
(\chirAss)^! \;\cong\; \chirAss \otimes \operatorname{sgn}.
\]
Consequently, $\Eone$--$\Eone$ Koszul duality exchanges $\Eone$-chiral
algebras with $\Eone$-chiral coalgebras, staying within the same
operadic type.
\end{proposition}

\begin{proof}
Classically, $\operatorname{Ass}^! \cong \operatorname{Ass} \otimes \operatorname{sgn}$
\cite[Theorem~7.4.1]{LV12}.  The chiral lift preserves Koszul duality:
$(\mathcal{P}^{\mathrm{ch}})^! \cong (\mathcal{P}^!)^{\mathrm{ch}}$
since the orthogonal complement commutes with the chiral tensor product
(computed fiber-by-fiber on $X^2 \setminus \Delta$, extended by $j_*$).
\end{proof}


\begin{definition}[$\Pinf$-chiral algebra = chiral Poisson]\label{def:pinf-chiral}
A \emph{$\Pinf$-chiral algebra} is an $\Einf$-chiral algebra equipped with a compatible
$L_\infty$ (homotopy Lie) structure.  This is the chiral analog of a Poisson algebra:
the commutative chiral product and the Lie bracket coexist via a Leibniz rule.
\end{definition}


\begin{remark}[$E_n$ terminology]\label{rem:En-terminology}
\index{$E_n$ terminology}
The subscripts in $\Einf$-chiral, $\Eone$-chiral, $\Pinf$-chiral refer
to the \emph{factorization level} (commutativity versus associativity of
the chiral product) inside the Beilinson--Drinfeld chiral tensor
category on a smooth algebraic curve.  Bare $E_n$ and $\En$ refer
instead to homotopy-theoretic little-disks algebras and topological
factorization on real $n$-manifolds \cite[\S5.1]{HA}.  In particular,
an $\Eone$-chiral algebra lives on a \emph{curve} (complex
dimension~$1$, real dimension~$2$) but is \emph{not} a topological
$E_1$-algebra: it is an associative algebra object in the chiral tensor
category of $\mathcal{D}_X$-modules.  When one compares the curve case
to the little-disks ladder, the relevant topological index is
$n = 2$, not $n = 1$.  The coincidence of notation is unfortunate but
standard in the Beilinson--Drinfeld literature, so the suffix
``-chiral'' is mathematically load-bearing throughout this monograph.
\end{remark}

The three types form a hierarchy
$\Einf\text{-chiral} \subset \Pinf\text{-chiral} \subset \Eone\text{-chiral}$.
Deformation quantization passes from $\Pinf$ to $\Eone$.
Bar-cobar duality operates at each level: $\Einf$ computes factorization homology; $\Eone$ connects to Yangians and quantum groups (Chapter~\ref{chap:yangians}).


\begin{definition}[Anomaly ratio]\label{def:anomaly-ratio}
\index{anomaly ratio|textbf}
For a Koszul chiral algebra $\cA$ with $c(\cA) \neq 0$, the \emph{anomaly ratio} is
\[
\varrho(\cA) := \frac{\kappa(\cA)}{c(\cA)}.
\]
For $\mathcal{W}$-algebras $\mathcal{W}^k(\mathfrak{g})$, this equals the exponent sum $\varrho(\mathfrak{g}) = \sum_{i=1}^r 1/(m_i+1)$, independent of the level $k$.  For Kac--Moody algebras $\widehat{\mathfrak{g}}_k$, $\varrho = (k+h^\vee)/(2k)$ depends on $k$.  The anomaly ratio measures the fraction of the conformal anomaly that contributes to the bar-complex curvature.
\end{definition}

\begin{definition}[Koszul conductor]\label{def:koszul-conductor}
\index{Koszul conductor|textbf}
For a chiral Koszul pair $(\cA, \cA^!)$, the \emph{Koszul conductor} is
\[
K(\cA) := c(\cA) + c(\cA^!).
\]
By Theorem~\ref{thm:central-charge-complementarity}, $K$ depends only on the root datum (not on the level $k$) for Kac--Moody and $\mathcal{W}$-algebra families.  The value $K = 2d$ for $\widehat{\mathfrak{g}}_k$ (where $d = \dim\mathfrak{g}$) is the Freudenthal--de Vries dimension; $K = 26$ for the Virasoro algebra is the bosonic string dimension; $K = 100$ for $\mathcal{W}_3$ is the $\mathcal{W}_3$-string dimension.
\end{definition}

\begin{definition}[Koszul spectrum]\label{def:koszul-spectrum}
\index{Koszul spectrum|textbf}
For a Koszul chiral algebra $\cA$, the \emph{Koszul spectrum} is the triple of invariants
\[
\operatorname{KSpec}(\cA) := (c(\cA),\; \kappa(\cA),\; \varrho(\cA))
\]
where $c$ is the central charge (Virasoro eigenvalue), $\kappa$ is the obstruction coefficient (Theorem~\ref{thm:genus-universality}), and $\varrho = \kappa/c$ is the anomaly ratio (Definition~\ref{def:anomaly-ratio}).  For $\mathcal{W}$-algebras, $\varrho = \varrho(\mathfrak{g}) = \sum 1/(m_i+1)$ is a root-system invariant (Corollary~\ref{cor:anomaly-ratio}).  The Koszul spectrum is additive under tensor products: $\operatorname{KSpec}(\cA \otimes \mathcal{B}) = \operatorname{KSpec}(\cA) + \operatorname{KSpec}(\mathcal{B})$ (for $c$ and $\kappa$; $\varrho$ is not additive).
\end{definition}

\begin{remark}[Three Koszul duality mechanisms]\label{rem:three-koszul-mechanisms}
Throughout this monograph, ``Koszul duality'' refers to three distinct but interrelated
mechanisms.  To prevent confusion, we distinguish them explicitly:
\begin{enumerate}[label=(\roman*)]
\item \emph{Francis--Gaitsgory chiral Koszul duality} \cite{FG12}:
an equivalence $\chirLie\text{-alg} \simeq \chirCom\text{-coalg}$ between chiral Lie algebras
and commutative chiral coalgebras, via the Chevalley/Primitives adjunction,
under pro-nilpotence hypotheses.  This operates at the level of operads
($\operatorname{Lie}^! = \operatorname{Com}$) in the chiral tensor category.

\item \emph{Bar--cobar resolution on the Koszul locus}: the bar and cobar
constructions exist for every augmented chiral algebra, but the counit
\[
\Omega(\bar B(\mathcal A)) \longrightarrow \mathcal A
\]
is a quasi-isomorphism only on the Koszul locus (Theorem~B), or in specific
completed/coderived regimes under additional hypotheses
(Appendix~\ref{app:nilpotent-completion}).  Off the locus, one retains
completed/coderived persistence
(Remark~\ref{rem:construction-vs-resolution}) but not ordinary
quasi-isomorphism.

\item \emph{Quadratic/Priddy duality} $\mathcal{A} \mapsto \mathcal{A}^!$:
for a quadratic (or more generally, augmented) chiral algebra, one constructs a
genuine \emph{Koszul dual algebra} $\mathcal{A}^!$.
When $\mathcal{A}$ is Koszul in the classical sense (the bar spectral sequence collapses
at $E_2$), the bar complex $\bar{B}(\mathcal{A})$ is quasi-isomorphic to $\mathcal{A}^!$.
For non-quadratic algebras, this requires nilpotent completion
(Appendix~\ref{app:nilpotent-completion}).
\end{enumerate}
Every theorem labeled ``Koszul duality'' specifies which of (i)--(iii) is invoked.
All three admit configuration space realizations via the Fulton--MacPherson compactification.
At higher genus, these three mechanisms fuse into a single
structure---the Feynman transform of the modular operad---where
the distinction between (i) and~(iii) becomes a filtration on the
bar complex by operadic weight, and the deformation-obstruction
complementarity of Theorem~C emerges as a Serre duality pairing
between the filtration layers.
\end{remark}


\begin{remark}[The master diagram: Verdier duality mediates]
\label{rem:master-diagram}
\index{master diagram}
\index{Verdier duality!as mediating mechanism}
Verdier duality on $\operatorname{Ran}(X)$ produces the Koszul dual.
The master diagram is the commutative square
\begin{equation}\label{eq:master-square}
\begin{tikzcd}[row sep=huge, column sep=huge]
\operatorname{Alg}_{\mathrm{aug}}(\operatorname{Fact}(X))
  \arrow[r, "\bar{B}_X"]
  \arrow[d, "(\cdot)^!"']
& \operatorname{CoAlg}_{\mathrm{conil}}(\operatorname{Fact}(X))
  \arrow[d, "\mathbb{D}_{\operatorname{Ran}}"]
\\
\operatorname{Alg}_{\mathrm{aug}}(\operatorname{Fact}(X))
  \arrow[r, "\bar{B}_X"']
& \operatorname{CoAlg}_{\mathrm{conil}}(\operatorname{Fact}(X))
\end{tikzcd}
\end{equation}
whose commutativity encodes Theorem~A:
$\mathbb{D}_{\operatorname{Ran}} \bar{B}_X(\cA) \simeq
\bar{B}_X(\cA^!)$.
The top row sends the algebra to its bar coalgebra; the right
column applies Verdier duality; the bottom row recognizes the result
as the bar of the Koszul dual.  The cobar functor
$\Omega_X$ is left adjoint to~$\bar{B}_X$, and the counit
$\Omega_X \bar{B}_X(\cA) \xrightarrow{\sim} \cA$
is Theorem~B.

At genus~$g \geq 1$, the square acquires a curvature parameter:
the conjectural universal Maurer--Cartan class
$\Theta_{\cA}$~\eqref{eq:universal-MC} would measure the failure of
strict commutativity, and the condition
$\kappa(\cA) = 0$ (anomaly cancellation) is precisely the
condition under which the square commutes in the derived
category rather than the coderived category.
This diagram recurs throughout the monograph: at genus~$0$ in
Chapter~\ref{chap:bar-cobar}, at higher genus in
Chapter~\ref{chap:higher-genus}, in the BV identification
of Chapter~\ref{ch:bv-brst}, and in every example of
Part~\textup{II}.
\end{remark}

\begin{remark}[Master functorial correspondence]
\label{rem:master-functorial}
\index{functorial correspondence}
Every row of the following table is a theorem.
\begin{center}
\small
\begin{tabular}{lll}
\textbf{Geometry} & \textbf{Algebra} & \textbf{Physics} \\ \hline
Verdier duality $\mathbb{D}_{\mathrm{Ran}}$
  & Koszul duality $\cA \leftrightarrow \cA^!$
  & S-duality / CPT \\[2pt]
Residue at collision divisor $D_{ij}$
  & Bar differential $d_{\mathrm{res}}$
  & OPE fusion \\[2pt]
FM boundary stratum
  & Tree contribution to $m_k$
  & Feynman diagram \\[2pt]
Arnold relation
  & $d^2 = 0$ (genus~$0$)
  & Factorization coherence \\[2pt]
Period integral on $\overline{\mathcal{M}}_g$
  & Obstruction class $\mathrm{obs}_g$
  & One-loop determinant \\[2pt]
Clutching map $\partial\overline{\mathcal{M}}_g$
  & Modular operad composition
  & String sewing \\[2pt]
DS reduction (BRST quotient)
  & Koszul functor $\mathrm{DS}$
  & Gauge fixing \\[2pt]
Log form $\eta_{ij} = d\log(z_i - z_j)$
  & Propagator in bar complex
  & Free-field propagator \\[2pt]
Forgetful $\overline{\mathcal{M}}_{g,n+1} \to
  \overline{\mathcal{M}}_{g,n}$
  & Forgetting a marked point
  & Integrating out a field \\[2pt]
Branch locus of dual Hilbert series
  & Discriminant $\Delta_{\cA}(x)$
  & Spectral curve invariant \\[2pt]
MC equation in $\barB(\widehat{\fg}_k)$
  & $L_\infty$ flatness
  & CS equations of motion \\[2pt]
$\int_{\overline{\mathcal{M}}_{0,n}} \barBgeom^{(0)}_n$
  & BRST cohomology pairing
  & Tree-level string amplitude
\end{tabular}
\end{center}
The first four rows are proved in Part~\textup{I}
(Chapters~\ref{chap:bar-cobar}--\ref{chap:higher-genus});
the next three are proved in
Part~\textup{II} (Chapters~\ref{chap:kac-moody-koszul}--\ref{chap:yangians});
the discriminant row is established as a formal invariant
in~\S\ref{sec:master-table}.
The last two rows are the new physics theorems:
the Chern--Simons identification
(Theorem~\ref{thm:cs-koszul-km}) and the genus-$0$ string
amplitude computation
(Corollary~\ref{cor:string-amplitude-genus0}).
The full physics column is made precise in
Chapter~\ref{ch:bv-brst} (BV-BRST identification) and
Chapter~\ref{ch:feynman} (Feynman diagram interpretation).
\end{remark}

\section{Main results with proof locations}
\label{sec:main-results-complete}

Chapter~\ref{ch:heisenberg-frame} demonstrates the four theorems
for the Heisenberg algebra; Theorem~A is proved in Chapter~\ref{chap:bar-cobar},
Theorems~B, C, D in Chapter~\ref{chap:higher-genus}.

\begin{theorem}[Central charge complementarity; \ClaimStatusProvedHere]\label{thm:central-charge-complementarity}
\index{central charge!complementarity|textbf}
\textup{[Regime: curved-central on the Koszul locus
\textup{(}Convention~\textup{\ref{conv:regime-tags})}.]}

For Koszul dual pairs related by the Feigin--Frenkel involution
$k \mapsto k' = -k - 2h^\vee$, the sum of central charges is
independent of the level:
\begin{enumerate}
\item[\textup{(a)}] \emph{Affine Kac--Moody.}
  $c(\widehat{\mathfrak{g}}_k) + c(\widehat{\mathfrak{g}}_{k'}) = 2 \dim \mathfrak{g}$.
\item[\textup{(b)}] \emph{Principal $\mathcal{W}$-algebras:}
  $c(\mathcal{W}^k(\mathfrak{g})) + c(\mathcal{W}^{k'}(\mathfrak{g}))
   = 2\operatorname{rank}(\mathfrak{g}) + 4h^\vee \dim \mathfrak{g}$.
\end{enumerate}
In particular, $\mathfrak{g} = \mathfrak{sl}_2$ in part~\textup{(b)} gives
$c + c' = 26$, recovering the bosonic string critical dimension from Koszul
duality alone.

For the $bc$--$\beta\gamma$ system, the analogous
complementarity $c + c' = 0$ follows from the Sym/Ext statistics exchange
(\S\ref{sec:fermion-boson-koszul}).  For the Heisenberg $\mathcal{H}_\kappa$,
the Koszul dual $\mathrm{Sym}^{\mathrm{ch}}(V^*)$ is curved, so its
central charge is not defined via the Sugawara construction; the
obstruction-level complementarity $\kappa + \kappa' = 0$ nonetheless holds.
\end{theorem}

\begin{proof}
\emph{Part (a).}  The Sugawara central charge is
$c(\widehat{\mathfrak{g}}_k) = k \dim\mathfrak{g}/(k + h^\vee)$.
Since $k' + h^\vee = -(k + h^\vee)$:
\[
c(k) + c(k') = \frac{k \cdot d}{k+h^\vee} + \frac{(k+2h^\vee) \cdot d}{k+h^\vee}
= \frac{2(k + h^\vee) \cdot d}{k+h^\vee} = 2d
\]
where $d = \dim\mathfrak{g}$.

\emph{Part (b).}  The Drinfeld--Sokolov central charge for the principal
$\mathcal{W}$-algebra is
\[
c(\mathcal{W}^k(\mathfrak{g})) = r - \frac{12|\rho|^2 (t-1)^2}{t},
\]
where $r = \operatorname{rank}(\mathfrak{g})$, $\rho$ is the Weyl vector, and
$t = k + h^\vee$.  This is the standard Drinfeld--Sokolov central charge,
obtained by applying the quantum Hamiltonian reduction of
\cite{FF} to the Sugawara tensor
(verified for $\mathfrak{sl}_2$ in
Theorem~\ref{thm:w-algebra-koszul-main} and for $\mathfrak{sl}_3$ in
Theorem~\ref{thm:w-koszul-precise}).
The involution sends $t \mapsto -t$, giving
$c(k') = r + 12|\rho|^2 (t+1)^2/t$.  Therefore:
\[
c(k) + c(k') = 2r + 12|\rho|^2 \cdot \frac{(t+1)^2 - (t-1)^2}{t}
= 2r + 48|\rho|^2.
\]
The Freudenthal--de~Vries identity
$|\rho|^2 = h^\vee \dim\mathfrak{g}/12$
(under the normalization $(\theta, \theta) = 2$ for the highest root~$\theta$; for non-simply-laced~$\mathfrak{g}$, this normalization matches the one used in the DS central charge formula above; see \cite{Kac}, Ch.~13)
yields $c + c' = 2r + 4h^\vee d$.

\emph{Verification.}
$\mathfrak{sl}_2$: $2(1) + 4(2)(3) = 26$.
$\mathfrak{sl}_3$: $2(2) + 4(3)(8) = 100$.
$\mathfrak{sl}_N$: $2(N{-}1) + 4N(N^2{-}1) = 2(N{-}1)(2N^2{+}2N{+}1)$.
\end{proof}

\begin{remark}[Relation to the bosonic string]\label{rem:bosonic-string-26}
\index{bosonic string!critical dimension}
The identity $c + c' = 26$ for $\mathfrak{sl}_2$ is the algebraic incarnation
of $c_{\mathrm{matter}} + c_{\mathrm{ghost}} = 0$: the ghost system is the
cobar of the matter Virasoro bar complex, and anomaly cancellation is
acyclicity of the bar-cobar adjunction (Theorem~\ref{thm:qme-bar-cobar}).
The general formula $2r + 4h^\vee d$ gives $100$ for $\mathfrak{sl}_3$;
this must not be confused with the quantum geometric Langlands duality
(Remark~\ref{rem:additional-dualities}).
\end{remark}

% [Arnold relations section removed — now demonstrated concretely
%  in Chapter~\ref{ch:heisenberg-frame}, \S3; proved in full generality
%  in Chapter~\ref{chap:config-spaces}, Theorem~\ref{thm:arnold-three};
%  three independent proofs in Appendix~\ref{app:arnold-relations}.]

% [Chiral Hochschild cohomology and periodicity theorems:
%  see Chapter~
ef{chap:hochschild-cyclic} for the full treatment,
%  including periodicity (Theorem~
ef{thm:periodicity-cohomology}).]

\section{Existence regimes for Koszul duals}

\begin{remark}[Strict, completed, and comparison regimes]
\label{rem:existence-regimes}\label{thm:existence-koszul}
The monograph does not use a single universal ``if and only if''
theorem for existence of Koszul duals.  It separates three regimes.
\begin{enumerate}[label=\textup{(\roman*)}]
\item \emph{Strict finite-type regime.}
  The constructions $\bar{B}_X(\cA)$ and $\Omega_X$ exist for every
  augmented chiral algebra, while inversion is theorematic only on the
  Koszul locus identified by the recognition data of
  Chapter~\ref{chap:koszul-pairs} and the higher-genus theorems of
  Chapter~\ref{chap:higher-genus}.
\item \emph{Completed frontier regime.}
  For weight-filtered or pronilpotent towers one seeks a completed bar
  object under explicit inverse-limit hypotheses.
  Appendix~\ref{app:nilpotent-completion} records that completed-model
  regime.  In the active MC4 frontier, however, the standard
  principal-stage completed M-level towers for $W_\infty$ and Yangian
  examples are already theorematic; what remains is the stronger
  H-level comparison package recorded in
  Chapter~\ref{chap:concordance}, namely the exact identities
  $\mathsf{C}^{\mathrm{res}}_{s,t;u;m,n}(N)
  =\mathsf{C}^{\mathrm{DS}}_{s,t;u;m,n}(N)$ and
  $\mathsf{K}^{\mathrm{line}}_{a,b}(N)
  =\mathsf{K}^{\mathrm{RTT}}_{a,b}(N)$ together with their finite
  detection on $\mathcal{I}_N$ and $\Delta_{a,0}(N)$.  It is not used
  to enlarge the proved core without those hypotheses.
\item \emph{Comparison regime.}
  Appendix~\ref{app:existence-criteria} compares the strict and
  completed regimes with the literature of quadratic and curved Koszul
  duality.  It is a map of paradigms, not a monograph-wide replacement
  for the family-specific theorem surface.
\end{enumerate}
The load-bearing distinction is therefore construction versus
resolution: $\bar{B}_X(\cA)$ may be constructed beyond the locus where
$\Omega_X\bar{B}_X(\cA)\simeq \cA$ is proved.
In particular, Appendix~\ref{app:nilpotent-completion} does not enlarge the
proved core unless its frontier hypotheses (finite generation, polynomial growth,
$E_2$-degeneration) are explicitly invoked.
\end{remark}


\section{Computations}

Part~II computes bar complexes for all major families
(Table~\ref{tab:master-invariants}).  The shared discriminant
$\Delta(x) = 1 - 2x - 3x^2$ (Theorem~\ref{thm:ds-bar-gf-discriminant})
and the universal generating function
$\sum_{g \geq 1} F_g x^{2g} = \kappa \cdot (\hat{A}\text{-genus})$
(Theorem~\ref{thm:universal-generating-function}) are consequences of
the proved scalar and spectral characteristic layers.

\section{Standing assumptions and conventions}
\label{sec:standing-assumptions}

\begin{convention}[Standing assumptions]\label{conv:standing-assumptions}
The following assumptions are in force throughout.
\begin{enumerate}[label=\textup{(SA\arabic*)}]
\item\label{SA:ground-field} The ground field is $\mathbb{C}$.  All vector spaces, algebras, and categories are over $\mathbb{C}$.

\item\label{SA:grading} We work in the cohomological grading convention: differentials have degree $+1$.  The bar construction uses the induced grading with suspension shifting degrees by $-1$.  See Appendix~\ref{app:sign-conventions} for detailed sign rules.

\item\label{SA:convergence} All infinite sums, products, and formal power series are understood in the appropriate completed tensor product or topological vector space.  When we write $\sum_{n \geq 0}$, convergence is in the formal (adic) topology unless stated otherwise.  Configuration space integrals converge absolutely on Fulton--MacPherson compactifications by Theorem~\ref{thm:FM-convergence}.

\item\label{SA:curves} Algebraic curves are smooth, projective, and geometrically connected over $\mathbb{C}$, unless stated otherwise.  ``Genus $g$'' always refers to the arithmetic genus.  Points on curves are closed points.

\item\label{SA:completion} Chiral algebras and vertex algebras are understood in their completed form when genus $g \geq 1$ constructions are involved.  The completion is with respect to the filtration induced by the order of poles at the diagonal; see Appendix~\ref{app:nilpotent-completion} for details.
\end{enumerate}
\end{convention}

\subsection{Ambient category and further conventions}
\label{subsec:ambient-category}
\label{subsec:grading-convention}
\label{subsec:augmentation-completion}
\label{subsec:coordinate-independence}
\label{rem:conventions-summary}

All $\mathcal{D}$-modules are holonomic.  The cobar
construction is the Verdier dual of the bar complex
(Definition~\ref{def:geom-cobar-intrinsic}), so
$d_{\mathrm{cobar}}^2 = 0$ follows from $d_{\mathrm{bar}}^2 = 0$
(Corollary~\ref{cor:cobar-nilpotence-verdier}).
Grading is cohomological ($|d| = +1$); see Appendix~\ref{app:signs}.
\index{grading convention!cohomological}
All chiral algebras are augmented
($\varepsilon\colon \mathcal{A} \to \omega_X$,
$\bar{\mathcal{A}} = \ker(\varepsilon)$);
\label{def:augmented-standing}
the bar complex is always reduced.
Forms $\eta_{ij} \in \Omega^1(\log D)$ and residue maps are intrinsic
to the FM compactification.
